% Created 2026-07-17 ven 21:36
% Intended LaTeX compiler: pdflatex
\documentclass[11pt]{article}
\usepackage[utf8]{inputenc}
\usepackage[T1]{fontenc}
\usepackage{graphicx}
\usepackage{longtable}
\usepackage{wrapfig}
\usepackage{rotating}
\usepackage[normalem]{ulem}
\usepackage{amsmath}
\usepackage{amssymb}
\usepackage{capt-of}
\usepackage{hyperref}
\usepackage[newfloat]{minted}
\author{marco robba}
\date{\today}
\title{Considerazioni Aggiuntive Al Testo D El Prof}
\hypersetup{
 pdfauthor={marco robba},
 pdftitle={Considerazioni Aggiuntive Al Testo D El Prof},
 pdfkeywords={},
 pdfsubject={},
 pdfcreator={Emacs 30.1 (Org mode 9.7.29)}, 
 pdflang={English}}
\begin{document}

\maketitle
\tableofcontents

\section{Per il DB}
\label{sec:org3e8b37b}
\subsection{Stazione}
\label{sec:orgf4dd051}
\subsubsection{Di tracciamento}
\label{sec:orgba31a23}
\begin{itemize}
\item Quali dati produce la stazione
\begin{enumerate}
\item Treno di passaggio
\begin{itemize}
\item dati dell'ingresso (timestamp e id treno)
\item dati dell'uscita  (timestamp e id treno)
\end{itemize}
sarebbero i sensori della stazione posti alle etrate e alle uscite della stazione
\begin{itemize}
\item Eccezione
un treno può entrarvi e non uscirvi capolinea
\begin{itemize}
\item gestire il caso per cui la stazione non è un capolinea ma comunque il treno lo si smonta e lo si ritira da quella stazione, molto spesso a causa di un guasto
\end{itemize}
un treno può uscirvi senza entrarvi
\end{itemize}
\end{enumerate}
\end{itemize}
\subsubsection{guasti}
\label{sec:org8b333c3}
\begin{itemize}
\item Stato del guasto(risolto/non risolto)
\begin{enumerate}
\item Gestione guasto nel canale informativo
Per non perdere i dati previo descritti allora è cosa buona immagazzianare tali per non perderli e reinviarli una volta che la connessione viene ristabilità e cancellarli una volta che saranno inviati tutti
\item Guasti vero e proprio della stazione
ne inibisce l'uso
\begin{itemize}
\item Esempi
\begin{itemize}
\item rottura di un binario
\begin{itemize}
\item per una particolare linea della stazione
\end{itemize}
\item rottura di un treno
\begin{itemize}
\item per una particolare linea della stazione
\begin{itemize}
\item treno esploso
\item treno deragliato (non credo sia necessario gestire questi casi)
\end{itemize}
\end{itemize}
\end{itemize}
\end{itemize}
\end{enumerate}
\end{itemize}
\subsection{Treno}
\label{sec:orga84df40}
\begin{itemize}
\item stato: fermo(per via di un guasto su una linea) o attivo guasto etc per gestire se lasciarlo fermo o farlo muovere
\begin{itemize}
\item le proprie informazioni
\item itinerario
\end{itemize}
\end{itemize}
\section{TAbelle Db}
\label{sec:org3dc95d2}
\subsection{Centrale}
\label{sec:org74465b4}
\begin{minted}[breaklines=true,breakanywhere=true,fontsize=\small]{sql}
CREATE TABLE
\end{minted}
\subsection{Stazione}
\label{sec:org70e383e}
\begin{minted}[breaklines=true,breakanywhere=true,fontsize=\small]{sql}
CREATE TABLE Stato (
    id_stazione                     INT PRIMARY KEY,
    nomeStazione                    VARCHAR(100) NOT NULL,
    regioneGeograficaDiAppartenenza VARCHAR(100),
    funzionanteONo                  BOOLEAN,
    flagFaultInCorso                BOOLEAN,
    timestampInizioUltimoFault      TIMESTAMP,
    timestampFineUltimoFault        TIMESTAMP
);

CREATE TABLE Tratte (
    id_tratta        INT PRIMARY KEY,  -- deve essere lo stesso del DB centrale
    stazionePartenza INT NOT NULL,
    stazioneArrivo   INT NOT NULL,
    FOREIGN KEY (stazionePartenza) REFERENCES Stato(id_stazione),
    FOREIGN KEY (stazioneArrivo)   REFERENCES Stato(id_stazione)
);

CREATE TABLE Guasto_Locale_Momentaneo (
    id_Guasto        INT PRIMARY KEY,
    statoRisoltoONo  BOOLEAN,
    tipo             VARCHAR(50),
    descrizione      VARCHAR(255),
    timestamp        TIMESTAMP
);

-- Tabella associativa per la relazione (0,N)-(0,N) Tratte <-> Guasto_Locale_Momentaneo
CREATE TABLE Tratta_Guasto (
    id_tratta INT NOT NULL,
    id_Guasto INT NOT NULL,
    PRIMARY KEY (id_tratta, id_Guasto),
    FOREIGN KEY (id_tratta) REFERENCES Tratte(id_tratta),
    FOREIGN KEY (id_Guasto) REFERENCES Guasto_Locale_Momentaneo(id_Guasto)
);

CREATE TABLE Treno (
    id_Treno      INT PRIMARY KEY,
    id_transito   INT,
    tempoUscita   TIMESTAMP,
    tempoPartenza TIMESTAMP
);

CREATE TABLE bufferChache (
    id_Evento INT PRIMARY KEY,
    tipo      VARCHAR(50),
    id_Guasto INT NULL,      -- (0,1) verso Guasto_Locale_Momentaneo
    id_Treno  INT NOT NULL,  -- (1,1) verso Treno
    FOREIGN KEY (id_Guasto) REFERENCES Guasto_Locale_Momentaneo(id_Guasto),
    FOREIGN KEY (id_Treno)  REFERENCES Treno(id_Treno)
);
\end{minted}
\end{document}
